% 图契约：本题的源、边界和推导逐项对应；箭头符号按图注解释。
\documentclass[tikz,border=7pt]{standalone}
\usepackage{fontspec,xeCJK,amsmath,amssymb}
\IfFontExistsTF{Times New Roman}{\setmainfont{Times New Roman}}{\setmainfont{TeX Gyre Termes}}
\IfFontExistsTF{Heiti SC}{\setCJKmainfont[AutoFakeBold=1.5]{Heiti SC}}{\setCJKmainfont[AutoFakeBold=1.5]{FandolHei-Regular}}
\usetikzlibrary{arrows.meta,bending,calc,patterns,positioning}
\definecolor{ink}{HTML}{243745}
\definecolor{blue}{HTML}{275DAB}
\definecolor{green}{HTML}{177C63}
\definecolor{red}{HTML}{BD4935}
\definecolor{orange}{HTML}{AA6C16}
\definecolor{purple}{HTML}{715393}
\tikzset{every node/.style={font=\fontsize{9.3}{12.5}\selectfont,text=ink},
axis/.style={-{Stealth[length=4pt]},line width=.65pt,draw=ink!65},
vec/.style={-{Stealth[length=5pt]},line width=1pt,draw=blue},
short/.style={-{Stealth[length=3pt]},line width=.8pt},
guide/.style={dashed,draw=ink!40,line width=.5pt},
card/.style={rounded corners=2pt,draw=blue!65,fill=blue!4,align=center,text width=6.5cm,inner xsep=9pt,inner ysep=7pt},
title/.style={font=\fontsize{10.5}{14}\selectfont\bfseries,text width=15cm,align=center},
label/.style={font=\fontsize{8.4}{11}\selectfont,align=center},
note/.style={label,text width=6.2cm}}
\begin{document}
\begin{tikzpicture}[x=1cm,y=1cm]
\path[use as bounding box] (-8,-3.5) rectangle(8,3.5);
\node[title] at (0,3.15) {M2-021\quad 偶极转矩形成小角度恢复};
\begin{scope}[shift={(-4,0)}]

\foreach\y in{-1.3,1.3}{\draw[vec](-2,\y)--(2,\y);}
\draw[ink,thick](-.8,-.6)--(.8,.6);\fill[blue](-.8,-.6)circle(2pt)node[left]{$-q$};\fill[red](.8,.6)circle(2pt)node[right]{$+q$};
\draw[vec,green](-.5,-.375)--(.5,.375)node[above left]{$\mathbf p$};
\draw[guide](0,0)--(1.2,0);\draw[orange](.6,0)arc(0:36.87:.6);\node[orange]at(.9,.25){$\beta$};
\node at(2,1.65){$\mathbf E$};
\node[note] at (0,-2.85) {稳定取向为偶极矩与电场同向，$pE>0$。};
\end{scope}
\node[card,draw=blue!65,fill=blue!4] at (3.9,1.8) {$I\ddot\beta=-pE\sin\beta$};
\node[card,draw=green!65,fill=green!4] at (3.9,0) {$\ddot\beta+(pE/I)\beta\simeq0$};
\node[card,draw=orange!65,fill=orange!4] at (3.9,-1.8) {$\omega=\sqrt{pE/I}$};
\end{tikzpicture}\end{document}